\section{The Metaphysics of Reintegration}

The fundamental aim of the \textbf{French Hermetic Tradition}, insofar as it is infused with \textbf{Martinism}, is \textbf{Reintegration}, that is, the restoration of the wholeness of man that was broken as a result of the Fall. In other words, it seeks the return to the Primordial State of Adam. Since the first task in the accomplishment of any task is the statement of its goal, it behooves us to try to understand more fully what that state would be like. Here we rely on the metaphysical teaching on the soul as well as the testimony of mystics and initiates.

A valuable resource that I am relying on here is the book \textit{Occult Phenomena in the Light of Theology} by the Trappist monk \textbf{Alois Wiesinger}. This work is a thorough and fair discussion of preternatural and occult phenomena including topics on psychoactive drugs, telepathy, clairvoyance, telekinesis, ghosts, possession, hypnosis, the subconscious, magic, astral travel, inter alia, which perhaps will deserve a separate review sometime. However, for our present purposes, we are more concerned with his teachings on the soul and its workings, on the inner experience of Adam, and the effects of the Fall. It is not necessary to adhere to a particular theology to gain from his discussion, but rather a commitment to the metaphysics of tradition and on personal experience.

First of all, Fr. Alois has to establish the structure of the soul. Here he relies principally on \textbf{St Thomas Aquinas}, although not in the abstract way of the philosophers, but as one who has understood the soul in his own interiority. In particular, as we have described here many times, there is the body, the vegetative, sensitive, and intellectual souls. However, the parts of the soul act together as a whole which, as such, is the form of the body. That is, the body, down to its cells and atoms, is shaped by the soul, or is its material manifestation. This teaching has consequences, which we cannot go into here; specifically, it means that manifestations such as one's sex, family, nation, etc., are not arbitrary and incidental, but rather are the reflection of something arising from the depths of one's spirit. This is the opposite of the modern view which would regard differences in talents, abilities, capacities, and so on, as cosmic injustices that require political and social interventions to rectify.

Along the way, Fr. Alois dismisses several other viewpoints. First of all, we can disregard materialism as incoherent and incompatible with common experience. Furthermore, he rejects

\begin{enumerate}
\item Platonic and Cartesian dualism with the two substances of body and soul,
\item Trichotomy, which then separates the soul and spirit,
\item Anthroposophy and Theosophy, which likewise identify even more substances with their various bodies.
\end{enumerate}

The fault of these conceptions is that they deny the integrality of the human being because of they regard the parts as separate substances. The Thomist (and we certainly don't restrict this solely to Thomas) view certainly incorporates the insights of these rejected views while accepting and even emphasizing the unity and indivisibility of the human substance. Were that not the case, then Reintegration would be impossible.

The soul, or psyche, thus described in ``body-bound". Hence, the way it knows is through the bodily senses and the images created in the mind. Fr. Alois then identifies another part that rises above the corporeal part of the soul, which he calls the anima spiritualis or ``spirit-soul" or nous. This spirit-soul can, in certain circumstances, withdraw from the body-bound part of the soul and ``allow its activity to reach out beyond the body. From this there result phenomena such as we encounter in occultism and to some extent in the mystic life."

There is a tendency in contemporary philosophers to restrict Thomas' teaching to the body-bound soul, which is certainly not his full teaching. This would reduce Thomism to effective materialism on the one hand; on the other hand, by denying the possibility of the spirit-soul, it falls into the Gnostic error of regarding material existence as a prison from which we need to escape for spiritual realization. It is true that ascetical practice, however, is necessary to free the spirit-soul from being totally body bound. Another way that the spirit-soul can perform its proper function is during the sleep state; this is a teaching of St Thomas and Fr Alois. This is why revelation or prophecy often comes in the form of dreams.

Fr Alois criticizes certain theologians who deny that there can be spiritual communication without the mediation of material sensations or images. To the contrary, God and angels do communicate with men directly; this is a dogmatic truth as well as the common experience of men from different traditions and cultures. Angels do communicate with men through thoughts; actually, this is the ordinary method, as angelic visions and materializations are the extraordinary way. Similarly, the spirit-soul, since it is effectively the equivalent to the angelic consciousness, can conceivably communicate with other spirit-souls through the medium of thought. This is actually the correct teaching of Thomas, since the soul is not completely submerged in the body, so it has possibilities that are not just those of the body.

\paragraph{Lost Tradition}
Fr Alois points to several specific precursors of the idea of the spirit-soul, not to mention the various unnamed seers, visionaries, and magicians. \textbf{Plato} in the \emph{Phaedrus} tells how men ``through divine madness become partakers of true prophecy," and mentions true dreams coming in the sleep state in the \emph{Republic}. He also points out that ``Aristotle knows of an exalted state of the soul in sleep, in which it withdraws into its own nature and has power over the future."

The Stoic \textbf{Posidonius of Apameia} wrote: ``there is another method of prophecy that proceeds from nature; this proves how great is the power of the spirit, when it has been released from the sensual organs of the body. This occurs especially in sleep and in ecstasy." Readers may recognize that the state of deep sleep is one of the states mentioned in the Vedanta, that of the causal body (i.e., the intellectual soul).

\textbf{Eudemos} also wrote on prophecy, claiming that the lower soul ``partakes of the divine in ecstasy and in dreams." The Delphic high priest Plutarch declares the daimonion to be ``the guardian spirit which, unlike the soul, is not completely united to the body, but reaches out beyond it and sometimes loosens its connection with it to wander abroad and communicate immediately with gods and spirits, whence it derives the gift of prophecy". Fr Alois claims the daimonion to be the spirit-soul. In Guenon's terms, this guardian angel is a higher state of man. Compare this to the episode in Acts 12:15 where the disciples believe that it couldn't be Peter at the door but ``it must be his angel"; they just took that possibility for granted.

He concludes with the Neoplatonists, especially \textbf{Plotinus}' attainment of the highest states four times. All accepted that, when transcending the body, the soul could function as a pure spirit, could contemplate God, apprehend truths prophecy, experience second sight, etc. This was the consistent teaching of antiquity and therefore the starting point of the Church fathers.

Unfortunately, that understanding was allowed to lapse into oblivion, and was replaced with a ``confused belief in demons". That is the situation today, where the occult phenomena described by Fr Alois are ignorantly and indiscriminately considered the products of the devil and are feared. However, this limited view arises from the lack of understanding of the powers of the spirit-soul.

Thus, attempts to revive such teachings tended to become marginalized, despite the positive writings of Thomas Aquinas (who recognized the power of clairvoyance) and other saints. Hence, Fr Alois then mentions figures such as \textbf{Swedenborg}, \textbf{Cornelius Agrippa} and \textbf{Paracelsus} who accepted certain occult powers as arising from the nature of man.

In our time, we see a small movement to reclaim these teachings as legitimate and necessary. The teachings of the soul are perfectly fine. Unfortunately, they have become desiccated, left to the philosophers, and divorced from the inner experiences that bring them to life and demonstrate their truth. On the other side, since the figures who have attempted to maintain that knowledge have been marginalized, and even condemned, such teachings have been set free from their traditional supports and often devolve into pseudo-mysticism and various aberrations.

That is why a Guenon or a Schuon insist on the attachment to a legitimate tradition, or why a Tomberg tries to bring Hermetism within the umbrella of the Tradition of the West, or a Mouravieff brings Fourth Way teachings back into Orthodoxy.


\paragraph{The Paradisal State}

\begin{quotex}
We know that it is very difficult to tell from a broken machine how its various parts are intended to operate. One can only learn that by seeing a sound machine in actual operation. The same is true of man, particularly when we are concerned with the most important part of him, namely his soul. In order to become acquainted with all its attributes and functions, it is necessary to study it in its sound condition … \flright{\textsc{Alois Wiesinger, OCSO}, \textit{Occult Phenomena}}

\end{quotex}
Fr Alois proposes to study the Adamic state in order to grasp the proper functioning of the soul. First, he takes care to define the natural and the supernatural.

\paragraph{The Natural}
The natural is what constitutes a substance or derives from it. Note that this is not what is commonly meant by ``nature" today, which is typically assumed to be how animals in the wild act. On the contrary, man has a specific nature that is more than his animal nature. In this sense, nature means:

\begin{enumerate}
\item What constitutes it.

Everything that inwardly constitutes the specific essence of a thing, both the essential and integrating parts of its being. 
\item What derives from it 

\begin{enumerate}
\item Everything that proceeds spontaneously: e.g., aptitudes, talents, and powers 
\item Everything that proceeds under the influence of another being. That is, things that must be taught such as an art, skill, or craft 
\end{enumerate}
\item What is required for it. I.e., everything necessary 

\begin{enumerate}
\item for its continued existence (nourishment, light, air) 
\item for its activity (the will for survival) 
\item for its development (education, society, government) 
\item for the attainment of its goal (knowledge of God, free will) 
\end{enumerate}
\end{enumerate}
These are called the ``demands of nature" as they are owed to man for them to exist as such. The supernatural and the preternatural are beyond nature and are not owed to man qua man. The former is all that transcends created nature and the latter whatever is created, but of a higher nature.

\paragraph{The Preternatural}
This includes everything of a higher nature. The example Fr Alois gives is this:

If a man makes an act of knowledge without the mediation of the senses and after the manner of the angels, then he transcends his own nature and is permitted to partake in the higher nature of the angel.

The type of knowing is called ``intuition", the direct seeing of higher things. This type of knowing is not based on sense images or the ``knowledge" of certain propositions. We can rephrase Fr Alois' description as man transcends the human state and reaches an angelic state, as Guenon would phrase it; they mean the same thing. There are two types of the preternatural:

\begin{enumerate}
\item According to the matter, or the thing done. This is a perfection to which we have no claim 
\item According to the form, the manner of receiving it. 
\end{enumerate}
The example above is of the first type, since the matter is from the angels, and man has no natural right to that. An example of the second type is infused knowledge; in this case, acquiring knowledge is natural, but having it infused is not. I would include prophecy or visions in this category.

\paragraph{The Supernatural}
The supernatural is beyond the natural in that it is not owed to man qua man. The supernatural is beyond all created things, including angels. Man is given a divine nature, beyond his human nature. Since this is freely given, and neither deserved nor earned by man qua man, it is often called a ``grace".

\paragraph{The Edenic State}
In the Edenic state, Adam had all that was proper to his human nature as described above. He also was given supernatural grace so he shared the same nature as God. In addition, he was given several preternatural privileges, in particular:

\begin{itemize}
\item Freedom from concupiscence 
\item Freedom from suffering and death 
\item The power of higher knowledge 
\item The faculties of pure spirits, i.e., his spirit-soul 
\end{itemize}
In this fully integral state, Adam shared a human nature, a divine nature, and an angelic nature. To grasp this state of being, we need to focus on the part of the soul that transcends the body. St Thomas describes three ways of knowing God:

\begin{enumerate}
\item After the Fall, we know God only in the mirror of his creatures. This was clearly described by Louis-Claude de Saint-Martin\footnote{\url{https://www.meditationsonthetarot.com/the-tableau-of-god}}
\item In Paradise, God was known via a spiritual light which he infused into the human spirit. This light was uncreated. This doctrine is still remembered in the Eastern churches as the Light of Tabor\footnote{\url{http://orthodoxwayoflife.blogspot.com/2010/02/truth-of-uncreated-light-of-god.html}}. 
\item In the beatific vision, God is known by the light of his glory 
\end{enumerate}
St Thomas explains that Adam's manner of knowing is similar to mystical contemplation and knows God as the angels do, at least in part. St Augustine wrote:

\begin{quotex}
Perhaps God spoke to the first human beings as he does to the angels, by illuminating their spirit with the unchanging light 

\end{quotex}
That is why ``reintegration" is often called ``illuminism".

\paragraph{Action follows Being}
So Adam had the powers of the corporal soul (tied to the body) and the spirit-soul (free of the body). The body was not a punishment, a prison, or a source of evil; rather, it gave Adam a new form of knowledge that was inaccessible to the spirit-soul. We should clarify Fr Alois' method at this point. On the one hand, he relies on metaphysical teachings and then draws all the necessary consequences from it. On the other hand, he studies empirical phenomena and works backwards to the kind of being that would produce such phenomena. Hence, by the analysis of mystical and occult phenomena, he is led to the conclusion that the spirit-soul is still active in men, even when dominated by the corporal soul. He sees these as residues of the primordial state and as clues to the attributes of that state.

That is justified on the metaphysical principle that ``action follows being", that is, whatever occurs is the result of the nature of the being. Hence, if Adam had angelic powers, then necessarily they required Adam to have a spirit-soul transcending the corporal soul.

\paragraph{The Causal State of Deep Sleep}

\textbf{Fr. Alois Wiesinger} makes a rather remarkable point about Genesis, when he writes:

\begin{quotex}
While he was creating woman, ``God cast a deep sleep over Adam", a sleep which in actual fact represented a great release from the senses. Theologians have been at some pains to explain the condition that is indicated by the word Tardemah. Though this word does not really mean ``ecstasy", which is the Septuagint rendering, it can nevertheless be rendered as an ecstatic sleep, that is to say, as a state of being in which the soul dwelt outside the world of sense and was active after the manner of pure spirits.

\end{quotex}
Obviously, this invites comparison with what \textbf{Rene Guenon} writes about the three human states in Man and his Becoming according to the Vedanta. He distinguishes:

\begin{enumerate}
\item The \emph{waking state}, corresponding to gross manifestation 
\item The \emph{dream state}, corresponding to subtle manifestation 
\item \emph{Deep sleep}, which is the causal and formless state 
\end{enumerate}
Moreover, he mentions the state of ecstatic trance, intermediate between deep sleep and death. Guenon, drawing inspiration from the Upanishads, repeats ideas from Genesis, even to the allegedly ``incorrect" translation of Tardemah as ``ecstasy". So it is worth the trouble to see what Guenon says about this state in Chapter XIV.

\begin{quotex}
These three are states of Atma, that is, the Person or Self, or, in other words, the possibilities of the Person. In the causal state, there is no longer individual existence, that is, the Person is transcendent. In the waking state, he is one corporeal being among many and in the dreaming state, he is creating an illusory world. In the causal state, the Person ``has become one", ``has identified himself with a synthetic whole of integral knowledge", etc. In this state, the Person is at the level of Sat, or pure Being. The intelligible Light is seized directly by intellectual intuition.

\end{quotex}
Thomas Aquinas also understands God as the act of Being and makes the same claim, albeit from a different frame of reference. He explains that Adam, in this state of deep sleep, knew God in His essence, i.e., at the level of Being.

This has several consequences. First of all, it seems to claim that God is known not only through his energies and but also in His essence; presumably, that is possible in the Primordial state. This is consistent with Thomas' understanding that to ``know" something is to participate in its essence. So to know is to be, so in a sense the knower is the thing known, at least in essence. So Adam in that state of deep sleep knew God in his essence and therefore ``was" God. However, he was not God in His existence, so we cannot say that Adam, as an individual being (i.e., in one of the states of manifestation), was God, but only as transcendent. With these ideas in mind, one should go back to Guido de Giorgio's descriptions of this state; it may make more sense. If anyone needs more, we can point to Johannes Tauler who explains:

\begin{quotex}
The more simple, pure, and naked in faith, so much the more praiseworthy, noble, and meritorious it is. This faith merits that God Himself in Himself, in His divine Essence, should be manifested to the soul in many wonderful ways.

\end{quotex}
\paragraph{Psychological Interpretation}
This causal state is beyond sensory imagery. In particular, therefore, it is also beyond any psychological process. Hence, the psychologist \textbf{Carl Jung} could not understand it. When he confronted certain yogis during his trip to India in 1938 about this state, he adamantly argued with them that it was simply unconsciousness, not a higher state of consciousness. He claimed, too, that the Indian mind was ``primitive" since he ``doesn't think" but instead ``perceives thought". That shows how alienated the modern mind is from its own past, since that attribute is closer to the medieval mind. \textbf{Ananda Coomraswamy} has pointed out that to understand the Indian mind, one need only refer to Western medieval and ancient writings.

Consider, for example, this description from \textbf{Saint Theresa d'Avila} about the difficulties that may assail the contemplative:

\begin{quotex}
There is a necessity of suffering the trouble of a Troop of Thoughts, importune Imaginations, and the impetuosities of natural Notions, not only, of the Soul through the dryness and disunion it hath, but of the Body also, occasioned by the want of submission to the Spirit, which it ought to have.

\end{quotex}
This is quite clear. The Saint is not ``thinking" those ``troops of thoughts" but instead is perceiving them as distractions to her own consciousness and arising from an outside source. This is not a matter of interpretation or an invitation to intellectual disputation. To perceive a thought, it is necessary to transcend the flow of thoughts in one's consciousness. Since life is short and art is long, a man or woman should start young to develop this skill.

\paragraph{Mystical Sleep}
Unlike the cultured despisers of their own Tradition who are so common in certain circles in our day, we can provide examples of this state in the true Tradition of the West. For example, \textbf{Fr. John Arintero} writes the following in \emph{The Mystical Evolution}:

\begin{quotex}
When faithful souls find themselves totally abandoned, arid, cold, mute, lacking affection, and unable to give vent to a single sigh, they raise their dim and darkened eyes to heaven and, directed by the obscure light of faith, they await mercy in silence. As they wait, sweetly enraptured and forgetful of all things, speaking not a word nor hearing nor seeing anything at all, they find themselves in a deep and profound silence which is at times changed into a mysterious sleep. Sometimes they pass entire hours in this condition, but to them it does not seem long because they experience a special attraction for that state.

\end{quotex}
\textbf{St. Francis de Sales} also writes about this mystical sleep that is beyond, and superior to, any conscious acts.

\begin{quotex}
When you find yourself possessed of this simple and pure filial trust in our Lord, remain in it, without attempting to perform any kind of conscious acts of intellect or will. This simple and trusting love, this loving sleep of the spirit in the arms of the Savior contains in itself as much excellence as you could ever hope to find. It is better to sleep on His sacred Breast than to be awake in any other place.

\end{quotex}
The crypto-modernists may be put off by such language, but if they are really traditionalists, they will understand how to relate such symbolic and theological language to metaphysical teachings.

\paragraph{Knowing the Nature of Things}

\textbf{Fr Alois Wiesinger} next provides some detail of the inner life of Adam, or, for our purposes, that of the primordial state. According to Fr Wiesinger, Adam ``intuitively beheld God, the creation of the world and the purpose thereof, the principles of law and morals and all that was necessary for him as head and instructor of the human race."

This is not purely speculation since it is based on metaphysical principles and also on the writings of saints and mystics. The latter, through their ascetical practice, were able to transcend the attachment to the corporal soul and reach similar states. \textbf{Saint Bernard of Clairvaux}, the patron of the Templars, explained it this way:

\begin{quotex}
It was only through sin that reason was thus imprisoned in the senses; once man also had a spiritual eye, that did not need the senses in order to know God, but this has now been clouded and darkened by sin and can only be cleansed for contemplation by asceticism.

\end{quotex}
This teaching, and it agrees with Oriental teachings insofar as they also seek to overcome the attachment to the senses by asceticism, has far reaching consequences. Specifically, those who have lost the primordial state and never made the effort of reintegration are ignorant of the purpose of the world and the principles of law and morals. At best, they can understand such things indirectly by studying the works of metaphysicians and contemplatives. Otherwise, this knowledge is lost and they can only grasp things through the senses.

In a traditional society, there will be those who know and the people accept this knowledge to the best of their ability, some better than others, and the result is a hierarchically arranged society. Obviously, there is no way to convince the modern mind of this because they regard it as an infringement on their intellectual freedom rather than as a liberation from sin and/or ignorance. That is acceptable and understandable. However, what is disappointing is the number of those who call themselves traditionalists who would reject such notions.

\paragraph{The Names of Things}
Adam's genius shows itself in the fact that he gave names to the animals. Saint Augustine rated it as an act of the highest wisdom and Pythagoras recognized that the wisest man was the one who first gave names to things. In the mentality of traditional man, the name indicates the nature of a thing, and the man who knows the name of a thing fundamentally knows its nature or essence. Here, without going into detail we can point to the Hindu doctrine of \emph{nama-rupa} (i.e., name and form) in which nama is the essential properties of a being.

There are two ways to grasp the nature of a thing. The ordinary way is by discerning the non-essential properties by abstracting from the sensual image; this requires protracted study and experience. The extraordinary way is by the intuitive understanding of pure spirits. However, in general, the mass of men fail to achieve much understanding at all, except perhaps in a particular field of study, because their understanding is hampered by passions, undisciplined imagination, evil disposition, the requirement to provide for one's support, and the weakness of forgetfulness.

We can postpone examples of all these for a future posting, but it should be clear that there can be no general agreement on anything because the knowledge of the true nature of things is so rare. As a matter of fact, ``naming" is by now the total opposite of the knowledge of the true nature of things. Instead, the modern mind believes it can actually change the nature of a thing by changing its name. I'm certain readers can find many examples of this as it goes by the term ``political correctness". This is nothing but group delusion and it is virtually an inescapable trap because there is no effort at all made to understand the true nature of things beyond the word for it.

\paragraph{The Law of Accidents}
If you know the true nature of things and circumstances, you will not make errors and mistakes. Hence, Adam was able to avoid suffering and all dangers to his health and well-being. Adam was not naturally immortal, but he had the possibility of not dying. Hence, he was not subject to accidents externally nor sickness and old age, internally. In this, even the saints and mystics cannot avoid. Some mystical thinkers such as Aurobindo Ghose and Charles Fillmore (the founder of the Unity churches) believed they could achieve immortality, but time proved them wrong. Transhumanists believe they can achieve immortality mechanically, biologically, or computationally, but it seems unlikely.

The will of Adam in the primordial state was mastery over the body and the lower parts of the soul, and it was not subject to concupiscence. Nevertheless, their bodies were powerful and healthy and they could still experience the delights of the senses. Of course, unlike men and women today whose inordinate and irrational desires lead them into all sorts of trouble, Adam easily avoided such problems.

Fr Wiesinger claims that Adam had powers over material things. Presumably he would have had paranormal powers, or ``siddhis", such as telekinesis, that are described in his book.

\flrightit{Posted on 2013-09-09 by Cologero }


\begin{center}* * *\end{center}

\begin{footnotesize}\begin{sffamily}

\texttt{jc00001 on 2013-09-09 at 08:01 said: }

\url{http://archive.org/details/occultphenomenai00wies} A free copy of the book I found online.


\hfill

\texttt{Michael on 2013-09-09 at 20:01 said: }

Another brilliant post. If the soul shapes the body, is it possible to alter the body through though alone?


\hfill

\texttt{Cologero on 2013-09-09 at 20:28 said: }

Obviously, Michael, there have been cases of mental healing. However, there is something even deeper and beyond thought, that is what is manifested in the body, so it can't be infinitely malleable. I take it to mean that you can actualize more of who you are, or, as Guenon puts it, manifest all your possibilities.

If you are referring to siddhis, or perhaps a ``resurrection body", I don't know how much is possible.


\hfill

\texttt{JA on 2013-09-10 at 14:39 said: }

so what then is the cause of body-soul conflicts, of the type described of Plotinus that he felt imprisoned in his body ?


\hfill

\texttt{Cologero on 2013-09-10 at 15:36 said: }

That is the result of the fallen state, JA, not so in the Primordial state.


\hfill

\texttt{JA on 2013-09-19 at 11:45 said: }

I am still having a hard time with the A-T seeing of form and essence as interconnected as opposed to Platonic matter-spirit dualism. There are some spiritual experiences I've seen that defy reason. 

To know a person's name is to know the person but how can one know a name ? Is your name the name your parents gave you for often in these days the most vulgar of reasons (oh how different from the Chinese who truly understand the importance of a name and spend years before giving their children a final name !) or something else more interior ? If my mother names me Lexus does that make me a car ?

From the Catholic perspective is the astral body the same gender as the physical body ? I think of Steiner's claim that the astral and mortal bodies are always opposite each other. A person I know who is male had a vision of himself in another dimension as a woman – when he then made contact with his angel, the angel told him that he should have a sex change to follow his true will.


\hfill

\texttt{scardanelli on 2013-09-19 at 12:43 said: }

``If my mother names me Lexus does that make me a car ?"

JA, I believe this is what Cologero refers to when he says ``the modern mind believes it can actually change the nature of a thing by changing its name." 

To truly know the name of something is to know it's essence, and one knows the essence of something through intuition, rather than through the senses or thinking mind. Tomberg refers to this type of knowing as gnosis. Perhaps the upcoming discussion of MOT may be fruitful in gaining this type of ``knowing."

It would seem to me that our given name simply refers to our ego, that is, the individuality rather than the personality, or true self. Thus, it is not a true name.


\hfill

\texttt{Cologero on 2013-09-19 at 12:47 said: }

JA, I think you profoundly missed the point about naming, probably because your dualism can't account for the connection between name and form. The ability to name things truly is a rarer power, as it required the recognition of name and form; I believe I pointed out that the modern mind thinks that just by naming things, it can change the essence. However, that is the result of the widespread philosophy of nominalism.

The angel's advice sounds suspect, so I would recommend to your friend that he try for two out of three.


\hfill

\texttt{Cologero on 2013-09-19 at 12:58 said: }

Thanks, Scardanelli. Good point about the personal name. There are many examples of men and women who have taken a new name that may be closer to their essence. Obviously, Sarah and Abraham. But even monks take a new name with their vows. E.g., Lao Tzu means ``old master", I believe, which better describes his essence.


\end{sffamily}\end{footnotesize}
